%% notes-tex/019-simb-multimodal/sections/1-summary.tex
%% SOURCE: every number here is carried from a later section and cited there. Nothing
%% originates in this section.

\section{Summary\secstatus{todo}}
\label{sec:summary}

\subsection{Terms used throughout\secstatus{todo}}
\label{sec:terms}

\begin{description}
\item[Head] The output module for one phenotype. The encoder and the perturbation operator
  are shared; a head is the last few layers that turn the shared representation into one
  phenotype's values. A run declares which heads are active, so \emph{betaxanthin head} and
  \emph{morphology head} name output modules, not models.
\item[Metabolome head] The head that predicts the 19 Mulleder amino-acid concentrations.
  It appears in two roles. As the \textbf{only} active head it is the amino-acid strand
  (\cref{sec:amino-acid}). Added \textbf{alongside} the betaxanthin head it is an
  \term{auxiliary head}: it is trained and scored, but the number being reported is still
  betaxanthin, and the question is whether sharing an encoder with it helps.
\item[Auxiliary head] Any head added to a run for its effect on a different head's score.
  Adding one changes two things at once, the gradient signal and the population of strains
  the run can use, which is why every comparison here pins the instance set.
\item[Strand] One phenotype prosecuted across a series of W\&B projects.
\item[\file{pearson_per_feature}] For each output feature (a reporter gene, a CalMorph
  parameter, one amino acid), the correlation across strains between prediction and
  measurement; then averaged over features. It goes to zero when a model predicts each
  feature's mean. The companion \file{pearson_per_instance} correlates across features
  within one strain and stays high under exactly that collapse, so it is never the
  objective here.
\item[\file{nmse}] Mean squared error divided by the variance of the target about each
  feature's mean, so $\texttt{nmse} = 1$ is exactly ``predict each gene's training mean''.
\item[\file{roll_max}] The scoring rule: the maximum of a centered five-point rolling mean
  of a validation curve. An upward-biased order statistic whose bias grows with the number
  of epochs run, so every score here travels with its epoch count.
\item[Ceiling] The highest correlation any predictor could reach against a target carrying
  measurement noise, estimated per strand from replicate structure.
\end{description}

\subsection{Where each strand stands\secstatus{todo}}

%% GENERATED by experiments/019-simb-multimodal/scripts/build_retrospective_tables.py
%% SOURCE: results/round_leaderboards.csv + the three ceiling JSONs. Do not hand-edit.
\begin{table}[htbp]
\centering\small
\begin{tabular}{lrrrrl}
\toprule
strand & ceiling & best & of ceiling & epochs (peak/last) & runs \\
\midrule
Expression (Kemmeren, Sameith) & 0.775 & \textbf{0.229} & 30\% & 4007/4105 & 163 \\
Expression, masked-label objective & 0.775 & \textbf{0.241} & 31\% & 9188/9997 & 16 \\
Morphology (CalMorph, Ohya) & 0.611 & \textbf{0.082} & 13\% & 27/65 & 23 (1 flat) \\
Expression and morphology, joint & -- & \textbf{0.051} & -- & 28/70 & 32 \\
Betaxanthin (Cachera) & 0.914 & \textbf{0.372} & 41\% & 254/999 & 44 \\
Betaxanthin with metabolome head & 0.914 & \textbf{0.430} & 47\% & 176/496 & 46 \\
Amino-acid pools (Mulleder) & -- & \textbf{0.211} & -- & 402/999 & 31 \\
Beta-carotene (Ozaydin) & 0.544 & \textbf{0.132} & 24\% & 799/999 & 39 \\
\bottomrule
\end{tabular}
\caption[]{Every strand, scored the same way. \emph{best} is \texttt{roll\_max}, the maximum of a centered five-point rolling mean of the validation curve, an upward-biased order statistic whose bias grows with epochs run, which is why the epoch column travels with it. \emph{epochs (peak/last)} gives where the best run peaked and where it stopped: a peak in the last tenth means the run was cut off while still improving. Ceilings are measured by three different estimators because the three measurement types require them; the amino-acid cell is empty because Mulleder has one replicate per strain, and no ceiling is estimable. \emph{flat} counts runs whose validation metric never left zero.}
\label{tab:strand-summary}
\end{table}


Read the epoch column with the score column. Expression and beta-carotene peak in the last
tenth of their best run, so those numbers are lower bounds set by the compute budget.
Morphology and the expression-morphology joint arm peak at epochs 27 and 28 of runs that
stopped at 65 and 70, which is not a result of any kind. Betaxanthin and amino acid peak in
the first half of a 999-epoch run and then overfit.

The two betaxanthin rows are not comparable to each other: the joint row trains on 3,832
strains against the plain row's 3,698, because the joint configuration restricts to
genotypes carrying both phenotypes. The within-grid paired comparison in
\cref{tab:bx-aa-paired} is the valid form of that contrast, and it comes out negative.

\subsection{Nine findings\secstatus{todo}}

\begin{enumerate}
\item \textbf{The expression strand contains no valid comparison between mechanisms.} All
  67 scored arms stopped at or below 300 epochs, against a validation curve that dips at
  200 to 300 and then climbs past its early peak. No peak has been observed at any budget
  up to 10,000 epochs. Every per-arm delta on record measures early-training transients.

\item \textbf{The expression objective is fighting itself, and it is the quantile loss
  rather than squared error.} On the longest run the quantile loss bottoms at epoch 463 and
  rises for the next 9,500 while Pearson climbs, but squared error bottoms at epoch 9,175,
  alongside the Pearson peak at 9,674; on the second-longest run the two are one epoch
  apart. \file{nmse} never returns below 1 after about epoch 400, so the model reaches
  $r = 0.236$ while being \emph{no better than predicting each gene's mean} in squared
  error. The cause is calibration, not capability: at the peak its predictions are
  $1.95\times$ more spread out than its own correlation justifies, and a post-hoc rescale
  by $r/s$ moves \file{nmse} from $1.010$ to $0.944$ while changing no correlation at all.
  Best-by-\file{mse} checkpointing would have picked within 500 epochs of the right model,
  so checkpointing is a narrower problem than previously recorded and is not what holds the
  strand at $0.24$.

\item \textbf{One expression run is a multi-day object, and this sets the campaign.} The
  best run took \textbf{91.4 hours, 3.81 days}, for 9,999 epochs at 32.9 s/epoch, and still
  peaked at $97\,\%$ of the way through. One GPU delivers $\approx 2{,}500$ epochs per day,
  so the number of arms a campaign can afford is fixed before any of them are chosen.

\item \textbf{The perturbation operator provably cannot express a pair term, and masked
  unmasking does not supply one.} At one perturbed gene the cross-attention softmax runs
  over a one-element key set, so the model reduces to $g(h_i + c_b)$ and nothing depends on
  which reporter is predicted given which gene was deleted. 16,200 of 32,760 attention
  parameters are dead. The scGPT-style scheme conditions on \emph{other genes' measured
  values}, and at $k = 0$, which is where scoring happens, every one of those features is
  exactly zero and the forward pass is identical to the unconditioned model. The two are
  different axes. Capacity is not the binding constraint either: a rank-32 reconstruction
  of the target reaches $0.727$. A rank-64 pair term already reaches $0.2274$ in 2.0 days
  against the masked objective's $0.2362$ in 3.8.

\item \textbf{Morphology has never been trained on its own data.} The Ohya screen holds
  4,718 deletions; every morphology run used the 1,440 that also carry expression, giving
  $n_{\mathrm{train}} = 1{,}161$. Its ceiling, measured here for the first time, is
  $\mathbf{0.611}$ over 278 features, with 201 of them individually above $0.5$. For a
  second-label experiment, fitness shares 4,220 strains with morphology against
  expression's 1,440.

\item \textbf{The free amino-acid profile predicts betaxanthin; the precursor does not.}
  Over the 4,432 deletions the two screens share, the 19-dimensional pool predicts
  betaxanthin at out-of-fold $r = \mathbf{0.298}$, while tyrosine alone, the chromophore
  precursor, predicts it at $\mathbf{0.064}$ and correlates marginally at $-0.076$. Roughly
  three quarters of the profile's power survives regressing out single-mutant fitness. On
  the 724 deletions with no Costanzo fitness record, which carry twice the betaxanthin
  spread, the profile predicts at $0.455$.

\item \textbf{The metabolome auxiliary head currently costs betaxanthin performance.}
  Paired within grid cell, joint minus control is $-0.0265 \pm 0.0159$ over the five
  comparable cells where the control learned the task, with 2 of 5 positive. The
  $+0.0203$ that motivated the replication reproduces, at $+0.0192$, in the one cell it came
  from. So the coupling exists in the data and the current head does not use it.

\item \textbf{The graph prior is at chance, and the mask deletes the one signal that
  exists.} Graph proximity does not predict which reporters respond to a deletion, on any of
  the nine graphs, in the orientation the model uses: AUC $0.4961$ to $0.5057$, at most
  $+0.0046$ over a degree-preserving control, and longer walks make it worse. Three of the
  nine graphs are also silent on most deletions, scoring only $29$ to $75\,\%$ of them. But
  \emph{direction} carries real signal: following tflink edges TF $\to$ target gives
  $0.5508$ and regulatory interaction $0.5239$, against controls near $0.501$. The mask
  builder symmetrizes, which averages that away. Nine attention heads are currently spent
  enforcing a target that does not predict response.

\item \textbf{Betaxanthin is the strongest case, and its strength is coverage.} yeast-GEM
  reaches $19\,\%$ of the screen and misses $73\,\%$ of its high producers. On genes a
  flux-based method has no features for, CGT's top 25 is enriched $4.4\times$ over base
  rate. Where both methods can be scored they are tied, they find largely different genes
  (9 shared of 29 each), and their union reaches 49 of 100.

\item \textbf{Configuration choice moves results more than method choice.} Betaxanthin grid
  cells span Spearman $0.013$ to $0.158$ against a between-method gap of $0.0014$ on AUC.
  The amino-acid strand's best run is $0.211$ and its median is $0.0498$.
\end{enumerate}

\subsection{The decision\secstatus{todo}}

\begin{itemize}
\item \textbf{Figure 6 is ready to be argued} on coverage rather than on accuracy, and the
  amino-acid material belongs in it as a coupling measurement, not as a joint-training
  result.
\item \textbf{Figure 3 is not ready, and its placeholder numbers are far from the measured
  ones.} A per-gene Pearson of $0.24$ against a ceiling of $0.775$ is not a headline. The
  cheapest path to a better one is not another architecture ladder: it is training
  morphology on its full dataset, finding where the expression curve actually turns, and
  checking whether the graph prior is the right prior at all.
\item \textbf{The graph-prior probe has run and it closed a planned arm.} The
  network-distance pair term is not worth building, the nine masked heads should be freed,
  and the mask should stop symmetrizing its two directed relations. Two cheap experiments
  remain unrun: morphology at full scale, and one expression run trained until the
  validation curve turns.
\end{itemize}
